\documentclass[exercice,a4]{thedocument}%
\usepackage{texmaths-tikz}%
\usepackage{pstricks}%
\usepackage{pst-bezier}%
\usepackage{pstricks-add}%
\usepackage{pst-eucl}%
\usepackage{pst-func}%
\usepackage{pst-tree}%
\usepackage{pst-circ}%
\usepackage{tikz}%
\usepackage{pgfplots}%
\usepackage{tkz-tab}%
\usepackage{tikzpagenodes}%
\usepackage{tabularray}%
\usepackage{tabularx}%
\usepackage{multirow}%
\usepackage{multicol}%
\usepackage{scratch3}%
\usepackage{enumitem}%
\usepackage[np]{numprint}%
\usepackage{pas-tableur}%
\usepackage{listings}%
\usepackage{graphicx}%
\usepackage{adjustbox}%
\usepackage{needspace}%
\usepackage{xcolor}%
\usepackage{eurosym}%
\usepackage{diagbox}%
\usepackage{colortbl}%
\usepackage{multido}%
\usepackage{icomma}%
\usepackage[french]{babel}%
\graphicspath{{../brevet/2026/}}%
\startdatakeys
    \source{26GENMATAA1}
    \BO{2026}
    \cycle{4}
    \competencesDuSocle{CA,RA,CO}
    \objectifsApprentissage{C4A1aa,C4A2ad,C4B2cb,C4D2aa,C4D3ac}
    \motsCles{angle,triangle équilatéral,carré,Scratch,programme,échelle}
\enddatakeys
\definecolor{gris}{gray}{0.5}%
\providecommand{\wideparen}[1]{\overset{\frown}{#1}}%
\begin{document}%
\providecommand{\degres}{^\circ}%
\begin{minipage}{0.48\linewidth}
	ABCD est un carré.

	ABF est un triangle équilatéral.

	AF $= 4$ cm.

	Les points F{}, A et E sont alignés.
\end{minipage}\hfill
\begin{minipage}{0.48\linewidth}
	\psset{unit=1cm}
	\begin{pspicture}(-3.5,-1.4)(3.5,4.8)
		%\psgrid
		\psline(-3.2,0)(3.2,0)
		\pspolygon(0;0)(3.2;120)(3.2;180)
		%\rput{30}(0,0){\psframe(0,0)(3.2,3.2)}
		\def\barb{\psline(-0.05,-0.1)(-0.05,0.1)\psline(0.05,-0.1)(0.05,0.1)}
		\rput{30}(0,0){\psline(0;0)(3.2;0)(4.525;45)(3.2;90)}
		\psarc*(0,0){0.6}{0}{30}
		\uput[d](1.6;180){4 cm}
		\uput[d](0,0){A}
		\uput[d](-3.2,0){F} \uput[d](3.2,0){E}\uput[ul](-1.6,2.772){B}\uput[r](3.2;30){D}\uput[ur](1.4,4){C}
		\rput(-1.6,0){\barb}\rput(1.6,0){\barb}
		\rput{30}(1.6;30){\barb}\rput{135}(2,3){\barb}\rput{120}(-0.8,1.4){\barb}
		\rput{600}(-2.4,1.4){\barb}\rput{30}(-0.2,3.6){\barb}
		\psdots[dotstyle=+,dotangle=45,dotscale=1.5](3.2,0)
		\rput(0,-1){\textit{Cette figure n'est pas en vraie grandeur.}}
	\end{pspicture}

\end{minipage}
\begin{enumerate}
	\item Justifier que la mesure de l'angle $\widehat{\text{EAD}}$ est égale à $30^\circ$.

	Dans la suite de l'exercice on utilisera l'échelle suivante : \textbf{10 pas} dans le programme représentent \textbf{1 cm} dans la réalité.

	\item Le \textbf{bloc triangle} ci-dessous permet de tracer un triangle équilatéral et le \textbf{bloc carré} permet de construire un carré :

	\begin{center}
		\begin{tblr}{colspec={|X[7cm,c]|X[7cm,c]|}}
			\hline
			\textbf{Le bloc triangle} & \textbf{Le bloc carré} \\
			\hline
			\begin{scratch}
				\initmoreblocks{définir \namemoreblocks{Triangle}}
				\blockrepeat{répéter \ovalnum{3} fois}
				{
					\blockmove{avancer de \ovalnum{\textbf{J}}}
					\blockmove{tourner \turnleft{} de \ovalnum{\textbf{K}} degrés}
				}
			\end{scratch}
			&
			\begin{scratch}
				\initmoreblocks{définir \namemoreblocks{Carré}}
				\blockrepeat{répéter \ovalnum{4} fois}
				{
					\blockmove{avancer de \ovalnum{\textbf{M}}}
					\blockmove{tourner \turnleft{} de \ovalnum{\textbf{N}} degrés}
				}
			\end{scratch}
			\\ \hline
		\end{tblr}
	\end{center}

	Donner les valeurs de \textbf{J}, \textbf{K}, \textbf{M} et \textbf{N} pour que les blocs triangle et carré permettent de construire un triangle équilatéral de 4 cm de côté et un carré de 4 cm de côté.

	\emph{Aucune justification n'est attendue.}

	\item Le \textbf{programme principal} utilise le bloc \textbf{Triangle} et le bloc \textbf{Carré}.

	L'instruction \og s'orienter à $90\degres$ \fg{} signifie que l'on s'oriente vers la droite.

	Écrire sur la copie le numéro de la figure obtenue grâce à ce programme.

	\emph{Aucune justification n'est attendue.}
\end{enumerate}

%PSTricks
\begin{tblr}{colspec={XXX}, vline{1,3,4}={solid}, vline{2}={3-4}{solid},
	cell{4}{1,2}={c},vspan=even}\hline
	\SetCell[r=1,c=2]{l}\textbf{Figure 1}&&\textbf{Programme principal}\\
	\SetCell[r=1,c=2]{c} %
	\psset{unit=0.8cm}
	\begin{pspicture}(0,-0.5)(8.5,1.5)
			\def\motif{\pspolygon[linecolor=blue](0;0)(0.75;120)(0.75;180)\psframe[linecolor=blue](0,0)(0.75,0.75)}
			\multido{\n=0.75+1.50}{6}{\rput(\n,0){\motif}}
		\end{pspicture}
	&& \SetCell[r=5,c=1]{c}
		\begin{scratch}[scale=0.85]
		\blockinit{quand \greenflag est cliqué}
		\blockmove{aller à x : \ovalnum{0} y : \ovalnum{-100}}
		\blockmove{s'orienter à ~\ovalnum{90}}
		\blockpen{effacer tout}
		\blockpen{stylo en position d'écriture}
		\blockrepeat{répéter \ovalnum{6} fois}
		{
			\blockmoreblocks{triangle}
			\blockmove{avancer de \ovalnum{50} pas}
			\blockmove{tourner \turnleft{} de \ovalnum{30} degrés}
			\blockmoreblocks{Carré}
			\blockmove{avancer de \ovalnum{50} pas}
			\blockmove{tourner \turnleft{} de \ovalnum{30} degrés}
		}
	\end{scratch} \\ \cline{1-3}
	\textbf{Figure 2}&\textbf{Figure 3}&\\
		{\psset{unit=0.7cm}
		\begin{pspicture}(-2,-2)(2,2)
			\def\dim{1}%%% modifiable
			\def\motif{\psline[linecolor=blue,linewidth=1.2pt](0,0)(\dim,0)(\dim,-\dim)(0,-\dim)(0,0)(\dim;60)(\dim,0)}
				\multido{\n=30+60,\na=0+60}{6}{\rput{\na}(\dim;\n){\motif}}
		\end{pspicture}}
		&
		{\psset{unit=0.9cm}
		\begin{pspicture}(-2,-2)(2,2)
			\def\motif{\psline[linecolor=blue,linewidth=1.2pt](0,0)(0,-1)(-0.8,-1)(-0.8,-0.2)(0,-0.2)(0,0)}
			\def\motiff{\psline[linecolor=blue,linewidth=1.2pt](0,0)(0.8;120)(0.8;180)(0,0)(1;120)}
			\multido{\n=15+60.0,\na=0+60}{6}{\rput{\na}(1.9;\n){\motif}}
			\multido{\n=15+60.0,\na=0+60}{6}{\rput{\na}(1.9;\n){\motiff}}
		\end{pspicture}} & \\ \cline{1-3}
	\SetCell[c=2]{l}\textbf{Figure 4}&&\\
	\SetCell[c=2]{c}
	{\psset{unit=0.8}
	\begin{pspicture}(0,-0.5)(9,1.5)
			\def\motif{\pspolygon[linecolor=blue](0;0)(0.75;120)(0.75;180)\psframe[linecolor=blue](0,0)(0.75,0.75)\psline(0.75,0)(1,0)}
			\multido{\n=1.00+1.75}{6}{\rput(\n,0){\motif}}
	\end{pspicture}}
	&&\\
	\hline
\end{tblr}
%Fin PSTricks
\medskip
\showthesource
\end{document}
